\chapter*{Phụ Lục}
\addcontentsline{toc}{chapter}{Phụ Lục}
\markboth{Phụ Lục}{Phụ Lục}

\section*{A. Ma Trận Sử Dụng Quyển XXXIII}

\begin{tongketchuong}
\textbf{Phần 1} phù hợp cho tiết sinh hoạt lớp, giáo dục kỹ năng sống, hoặc các buổi nói chuyện về quản trị cảm xúc và trách nhiệm cá nhân.

\textbf{Phần 2} phù hợp cho thảo luận về văn hóa lớp học, kỷ luật, phản biện, và cách người lớn giữ uy tín mà không tự thần thánh hóa mình.

\textbf{Phần 3} dùng tốt trong họp phụ huynh, tư vấn chủ nhiệm, hoặc trao đổi ba bên khi học sinh sa sút.

\textbf{Phần 4} phù hợp cho hướng nghiệp, giáo dục thái độ lao động, và những buổi nói chuyện về ảo tưởng thành công thời mạng xã hội.
\end{tongketchuong}

\section*{B. Gợi Ý Triển Khai 12 Chương}

\begin{goigiaovien}
\textbf{Chương 1--3:} mỗi chương có thể mở bằng một câu nói rất quen của học sinh rồi cho cả lớp đoán xem câu đó là thật lòng, chống chế, hay nửa nọ nửa kia.

\textbf{Chương 4--6:} nên dạy trong không khí bớt phòng thủ. Mục tiêu không phải bắt lỗi thầy cô, mà tập cho học sinh hiểu rằng quyền lực càng lớn càng phải đi cùng khả năng tự sửa.

\textbf{Chương 7--9:} nếu dùng với phụ huynh, nên bỏ giọng châm biếm và giữ phần cốt lõi: nghe sự thật về con, ngừng so sánh, và chuyển từ đổ lỗi sang kế hoạch.

\textbf{Chương 10--12:} nên kết hợp với câu chuyện thật từ cựu học sinh, người đi làm sớm, hoặc người từng đổi ngành để học sinh thấy khoảng cách giữa clip ngắn và đời thật.
\end{goigiaovien}

\section*{C. 24 Câu Hỏi Gợi Mở Trên Lớp}

\begin{multicols}{2}
\begin{enumerate}[leftmargin=1.2cm]
  \item Khi nào một lời than mệt là thật, và khi nào nó chỉ là cách né việc khó?
  \item Vì sao những từ đẹp như chữa lành, chữa trị, yêu bản thân lại dễ bị dùng sai?
  \item Em từng làm gì chỉ vì sợ bị tách nhóm?
  \item Người lớn xin lỗi có làm họ mất uy không?
  \item Nội quy nào ở trường cần được giải thích lại tử tế hơn?
  \item Khác biệt giữa phản biện và cãi cho thắng là gì?
  \item Bênh con đến đâu thì còn là tình thương, và từ đâu thì thành hại con?
  \item Em còn mang trong đầu câu so sánh nào từ gia đình?
  \item Một buổi họp phụ huynh tốt nên kết thúc bằng điều gì?
  \item Vì sao mạng xã hội hay kể đoạn hào quang mà không kể đoạn nền?
  \item Thế nào là trải nghiệm làm mình lớn lên thật?
  \item Tự do mà không có kỹ năng thì dễ biến thành gì?
  \item Không phải mọi áp lực đều là độc hại. Em đồng ý không?
  \item Khi bị góp ý thẳng, làm sao phân biệt bị xúc phạm với bị chỉnh cho tốt hơn?
  \item Sĩ diện đóng vai trò gì trong phần lớn các cuộc cãi ở trường?
  \item Vì sao người lớn nhiều khi đúng lý mà vẫn không thuyết phục được học sinh?
  \item Nếu được viết lại một nội quy, em sẽ sửa gì trước tiên?
  \item Điều gì khiến phụ huynh sợ nhất khi nghe thầy cô phản ánh về con?
  \item Một học sinh sa sút thường cần thêm kỷ luật, thêm hỗ trợ, hay cả hai?
  \item Tại sao nhiều bạn trẻ thích làm lớn nhưng ngại phần lặp lại và âm thầm?
  \item Đi làm sớm giúp người trẻ trưởng thành ở điểm nào?
  \item Có phải cứ mềm mỏng mới là thấu hiểu?
  \item Có phải cứ nghiêm mới là có trách nhiệm?
  \item Sau khi đọc xong quyển này, phe nào em hiểu hơn trước?
\end{enumerate}
\end{multicols}

\section*{D. Phiếu Ghi Chú Sau Mỗi Buổi Dạy}

\newcommand{\phieughi}{%
  \begin{tcolorbox}[
    colback=white,
    colframe=booksecondary!45,
    arc=4pt,
    boxrule=0.6pt,
    height=5.2cm,
    valign=top
  ]
  \textit{Ngày dạy / lớp: \hrulefill}

  \vspace{0.25cm}
  \textit{Chương sử dụng: \hrulefill}

  \vspace{0.25cm}
  \textit{Phản ứng nổi bật của học sinh:}

  \vspace{0.25cm}
  \hrulefill

  \vspace{0.25cm}
  \hrulefill

  \vspace{0.25cm}
  \textit{Điều cần điều chỉnh cho lần sau:}

  \vspace{0.25cm}
  \hrulefill
  \end{tcolorbox}
  \vspace{0.35cm}
}

\phieughi
\phieughi
\phieughi
\phieughi
\phieughi
\phieughi

\section*{F. Kịch Bản Sinh Hoạt 4 Tuần}

\begin{goigiaovien}
\textbf{Tuần 1: Gọi Đúng Tên Áp Lực}

Mở bằng Chương 1 hoặc Chương 2. Cho học sinh viết ẩn danh một câu các em thường dùng khi mệt hoặc không muốn làm việc. Giáo viên chọn vài câu tiêu biểu, cả lớp phân tích câu nào là cần giúp thật, câu nào là cách né tránh, câu nào là cả hai.

\textbf{Tuần 2: Phản Biện Và Kỷ Luật}

Dùng Chương 4, 5 hoặc 6. Chia lớp thành hai nhóm: nhóm bảo vệ uy tín của người lớn, nhóm bảo vệ quyền được phản biện của học sinh. Sau đó đổi vai để mỗi bên phải tự nói lý cho phía còn lại.

\textbf{Tuần 3: Nhà Trường Và Gia Đình}

Dùng Chương 7, 8, 9. Mời học sinh ghi lại một câu của cha mẹ khiến các em tổn thương nhất, nhưng không nêu tên phụ huynh. Giáo viên tổng hợp thành các mẫu câu quen thuộc rồi dẫn cả lớp nhìn xem phía sau các câu đó là lo lắng, áp đặt hay so sánh.

\textbf{Tuần 4: Tương Lai Không Chỉ Có Hào Quang}

Dùng Chương 10, 11, 12. Mỗi nhóm chọn một ảo tưởng quen thuộc: bỏ học vẫn giàu, trải nghiệm là đủ, đi làm sớm là trưởng thành sớm. Sau đó nhóm phải chỉ ra mặt sáng, mặt tối và điều kiện để lựa chọn ấy không trở thành trả giá mù quáng.
\end{goigiaovien}

\section*{G. 12 Mẫu Câu Mở Đầu Tiết Thảo Luận}

\begin{multicols}{2}
\begin{enumerate}[leftmargin=1.2cm]
  \item Có câu nào em nói quen miệng nhưng thật ra chỉ để tự cứu sĩ diện không?
  \item Một người lớn xin lỗi trước lớp làm em nể hơn hay coi thường hơn?
  \item Vì sao bị so sánh lâu ngày khiến người ta lười cố hơn là muốn cố?
  \item Có phải mọi nội quy gây khó chịu đều vô lý?
  \item Điều gì khiến học sinh hay dùng từ ``áp lực'' hơn là nói rõ vấn đề thật?
  \item Nếu em là phụ huynh của chính mình, em sẽ lo điều gì nhất?
  \item Tại sao nhiều bạn thích làm điều lớn lao nhưng chán phần bền bỉ nhỏ mỗi ngày?
  \item Một lời góp ý tệ là vì nội dung hay vì cách nói?
  \item Người lớn có cần học lại cách nói chuyện với Gen Z không?
  \item Học sinh có đang đòi được hiểu nhiều hơn mức mình chịu tự hiểu người khác không?
  \item Trải nghiệm nào làm em trưởng thành thật, không chỉ có ảnh đẹp?
  \item Sau một cuộc cãi nhau trong lớp, ai nên xuống giọng trước?
\end{enumerate}
\end{multicols}

\section*{H. Mẫu Cam Kết Sau Khi Đọc Sách}

\begin{tcolorbox}[
  colback=booksecondary!6,
  colframe=booksecondary!70!black,
  arc=4pt,
  boxrule=0.8pt,
  title={Cam Kết Cá Nhân \quad\textit{(Personal Commitment)}}
]
\textit{Sau khi đọc Quyển XXXIII, tôi muốn chỉnh lại 3 điều trong cách học, cách dạy, hoặc cách đồng hành của mình:}

\vspace{0.3cm}
1. \hrulefill

\vspace{0.55cm}
2. \hrulefill

\vspace{0.55cm}
3. \hrulefill

\vspace{0.45cm}
\textit{Người chứng kiến / người nhắc tôi giữ lời: \hrulefill}
\end{tcolorbox}

\section*{I. Lời Kết Nhỏ}

\begin{center}
{\itshape\color{bookprimary}``Lớp học thời nào cũng có lý sự, nhưng thời nào cũng cần người chịu ngồi lại để nghe nhau tới cùng.''}
\end{center}

\ngancach